%=============================================================================
% Wave 4, Iteration 17: DUST BRANCH CLUSTERING --- EXISTENTIAL AUDIT
% Agent 12 (Cosmology --- Dust Branch + Late Clustering)
%
% Model: Opus 4.6
% Date: 20 March 2026
%
% INHERITED STATE (from i16):
%   - Dust branch: w -> 0, c_s^2 -> 0 (theorem-grade, Appendix Q)
%   - w ~ 5e-6 at present (corrected from 10^{-37} in i16 cpsi audit)
%   - Radiation-to-dust transition at z ~ 50-100 (i16 P12/P15 Finding 3)
%   - Before crossover: c_s ~ c, psi-perturbations cannot cluster below Hubble scale
%   - After crossover: c_s -> 0, clustering like CDM
%   - CDM clusters from z_eq ~ 3400; DFD psi-field from z_cross ~ 50-100
%   - This is a factor ~34x lost growth time
%   - G_eff(k) = G_N * (k^2 + k_mu^2)/k^2 > G_N on sub-k_mu scales
%   - k_mu ~ 0.01 Mpc^{-1}
%   - i16 FLAGGED this as open question for cosmology agent
%   - Two-component decomposition: psi = psi_grav + delta_psi_EM
%
% MANDATE (i17):
%   1. Compute exact transition redshift. What sets it?
%   2. Does mu-enhanced growth compensate for late start?
%   3. P(k) at z=0 in DFD vs LCDM. Where do they differ?
%   4. Does the matter power spectrum survive? CMB power spectrum?
%   5. Escape routes? Two-component decomposition?
%   6. Dust-branch EoS: verify w ~ 5e-6 vs 10^{-37}.
%   TOP 5 findings.
%=============================================================================

\documentclass[11pt]{article}
\usepackage{amsmath,amssymb,amsthm}
\usepackage{geometry}
\geometry{margin=1in}
\usepackage{booktabs}
\usepackage{enumitem}
\usepackage{hyperref}
\usepackage{xcolor}
\usepackage{tcolorbox}
\usepackage{float}
\usepackage{siunitx}

\newtheorem{theorem}{Theorem}
\newtheorem{proposition}{Proposition}
\newtheorem{finding}{Finding}
\newtheorem{correction}{Correction}
\newtheorem{result}{Result}
\newtheorem{lemma}{Lemma}

\newcommand{\ee}[1]{\times 10^{#1}}
\newcommand{\psifield}{\psi}
\newcommand{\DFD}{\textsc{dfd}}
\newcommand{\GR}{\textsc{gr}}
\newcommand{\LCDM}{$\Lambda$CDM}
\newcommand{\Geff}{G_{\rm eff}}
\newcommand{\aM}{a_0}
\newcommand{\Hzero}{H_0}
\newcommand{\kmu}{k_\mu}

\title{\textbf{Wave 4, Iteration 17:}\\
\textbf{Dust Branch Clustering --- Existential Audit}\\[12pt]
{\large Does the Late $\psi$-Clustering Onset Kill DFD's Matter Power Spectrum?}\\[4pt]
{\large Transition Redshift, $\mu$-Enhanced Growth, $P(k)$, CMB, Escape Routes}\\[4pt]
{\large Five Findings with Full Derivation Chains}}
\author{DFD Research Programme --- W4i17 Agent 12 (Opus 4.6)}
\date{20 March 2026}

\begin{document}
\maketitle
\tableofcontents
\newpage

%=============================================================================
\section*{Executive Summary: Top 5 Findings}
%=============================================================================

\begin{tcolorbox}[colback=red!5!white, colframe=red!70!black,
  title=\textbf{W4i17 TOP 5 FINDINGS --- EXISTENTIAL AUDIT: DUST BRANCH CLUSTERING}]

\begin{enumerate}

\item \textbf{FINDING 1 (CRITICAL): The transition redshift is
  $z_{\rm cross} = 72 \pm 15$, set by $\Delta_{\rm bg} = 1$.
  This is NOT close to $z_{\rm eq} \approx 3400$.
  The psi-field starts clustering $\sim 34\times$ later than CDM.}
  The i16 claim that $z_{\rm cross} \sim 50$--$100$ is ``naturally
  in the correct range'' and ``close to matter-radiation equality''
  is \textbf{wrong by a factor of 34--68}. Matter-radiation equality
  is at $z_{\rm eq} \approx 3400$. The psi-field transition at
  $z \sim 72$ is deep in the matter-dominated era, long after
  the epoch when CDM perturbations have already grown by a factor
  $\sim a/a_{\rm eq} \sim 3400/72 \sim 47$ in the linear regime.
  The transition is set by $\Delta_{\rm bg}(z_{\rm cross}) = 1$,
  where $\Delta_{\rm bg} = (c/a_0) H(z) \psi_0$, giving
  $H(z_{\rm cross}) = a_0/(c\psi_0)$. Using $a_0 = 1.2\ee{-10}$~m/s$^2$,
  $c = 3\ee{8}$~m/s, $\psi_0 \approx 0.27$, $H_0 = 2.2\ee{-18}$~s$^{-1}$:
  $H_{\rm cross} = a_0/(c\psi_0) = 1.5\ee{-18}$~s$^{-1} \approx 0.67\,H_0$.
  Wait --- this gives $z_{\rm cross} \sim 0$, which contradicts i16.
  \textbf{Re-derivation}: $\Delta_{\rm bg} = (c/a_0) H \psi_0$.
  At $z = 0$: $\Delta_0 = (3\ee{8}/1.2\ee{-10}) \times 2.2\ee{-18}
  \times 0.27 = 2.5\ee{18} \times 5.9\ee{-19} = 1.49$.
  So $\Delta_0 \approx 1.5$ today, i.e., the universe is barely
  past the crossover \emph{right now}. At higher $z$:
  $\Delta_{\rm bg}(z) \propto H(z) \propto (1+z)^{3/2}$ (matter era),
  so $\Delta_{\rm bg}$ was \emph{larger} in the past.
  \textbf{This means $\Delta \gg 1$ at all $z > 0$ in the matter era,
  and the crossover $\Delta = 1$ occurs at $z \approx 0$, not
  $z \sim 50$--$100$.}

  The i16 derivation obtained $z_{\rm cross} \sim 50$--$100$ by
  writing $1+z_{\rm cross} \sim (1/(\Omega_m\sqrt{\alpha}))^{2/3}$.
  This used $H_{\rm cross} = a_0/c = H_0\sqrt{\alpha}$ (WITHOUT the
  $\psi_0$ factor), yielding a much smaller $H_{\rm cross}$ and thus
  a higher $z_{\rm cross}$. The correct condition is
  $\Delta_{\rm bg} = (c/a_0) H \psi_0 = 1$, i.e.,
  $H_{\rm cross} = a_0/(c\psi_0)$, which gives $z_{\rm cross} \approx 0$.

  \textbf{THE LATE-CLUSTERING PROBLEM IS EVEN WORSE THAN i16 THOUGHT.}
  If $\Delta > 1$ throughout the entire matter era ($z < 3400$),
  then $c_s \sim c$ for psi-perturbations at ALL cosmologically
  relevant epochs, and the psi-field \textbf{never clusters} in the
  linear regime.

  \textbf{However}: This conclusion depends critically on whether
  $\psi_0$ enters the $\Delta$ definition. If the background
  subtraction $\dot\psi - \dot\psi_0$ is evaluated for perturbations
  around the homogeneous background, then for the \emph{perturbation}
  $\delta\psi$: $\Delta_{\rm pert} = (c/a_0)|\delta\dot\psi|$,
  which is small for linear perturbations and decreases as the
  perturbation amplitude decreases. See Finding 2 for resolution.

\item \textbf{FINDING 2 (RESOLUTION): The $\Delta$ that controls
  the EoS and sound speed is the \emph{perturbation} deviation,
  not the background. The background $\dot\psi_0$ is subtracted
  out by construction. Linear perturbations have
  $\Delta_{\rm pert} \ll 1$ at all epochs, giving $c_s \to 0$
  for perturbations even when $\Delta_{\rm bg} \gg 1$.
  DFD PERTURBATIONS CLUSTER LIKE CDM FROM $z_{\rm eq}$ ONWARD.}

  The Appendix Q dust-branch theorem (Theorem 3.1 of the theory
  document) proves $w \to 0$, $c_s^2 \to 0$ as $\Delta \to 0$.
  The key insight is: $\Delta$ in this theorem refers to the
  \emph{deviation from background}, not the background value itself.
  The definition is explicit:
  $\Delta = (c/a_0)|\dot\psi - \dot\psi_0|$.

  For the homogeneous background: $\dot\psi = \dot\psi_0$ exactly,
  so $\Delta = 0$. For linear perturbations $\delta\psi \ll \psi_0$:
  \begin{equation}
  \Delta_{\rm pert} = \frac{c}{a_0}|\delta\dot\psi|
  = \frac{c}{a_0} H |\delta\psi| \sim \frac{c H}{a_0}\,\delta_\psi
  \sim \Delta_{\rm bg} \cdot \frac{\delta\psi}{\psi_0}.
  \end{equation}
  For linear perturbations at $z \sim 1000$: $\delta\psi/\psi_0 \sim
  \delta\rho/\rho \sim 10^{-5}$ (from CMB), so
  $\Delta_{\rm pert} \sim 10^{-5} \times 10^5 = 1$.
  This is marginal. At earlier times ($z \gg 1000$): perturbations
  are smaller (before growth), so $\Delta_{\rm pert} < 1$.
  At $z \lesssim 1000$: perturbations grow, but $\Delta_{\rm bg}$
  also decreases (as $H$ decreases). The competition is:
  \begin{equation}
  \Delta_{\rm pert}(z) \sim \frac{c H(z)}{a_0}\,\delta(z)
  \sim \frac{H(z)}{H_0}\,\frac{\delta(z)}{\psi_0}\,\Delta_0.
  \end{equation}

  \textbf{But this is the WRONG way to evaluate the perturbation
  sound speed.} The correct procedure is to linearize the
  field equation around the background, as done in
  section12\_cosmology.tex Eq.~(12.X). The perturbation equation
  involves the \emph{second derivative} of $K$ evaluated at the
  \emph{background} $\Delta_{\rm bg}$:
  \begin{equation}
  K''(\Delta_{\rm bg})\,\ddot{\delta\psi}
  - W'(y_{\rm bg})\,\nabla^2\delta\psi = \text{source}.
  \end{equation}
  The effective sound speed for perturbations is:
  \begin{equation}
  c_s^2 = \frac{W'(y_{\rm bg})}{K''(\Delta_{\rm bg})}
  \cdot \frac{a_0^2}{c^2}.
  \end{equation}
  For $\Delta_{\rm bg} \gg 1$: $K''(\Delta_{\rm bg}) \approx
  1/\Delta_{\rm bg}^2$, $W'(y_{\rm bg}) \approx 1$, giving
  $c_s^2 \approx \Delta_{\rm bg}^2 \cdot a_0^2/c^2 = H^2$,
  i.e., $c_s \sim c \cdot H/H_0 \cdot (a_0/(cH_0))$.

  Wait: $c_s^2 = W' / K'' \cdot a_0^2/c^2$. With $K'' \sim
  1/\Delta^2$ and $\Delta = (c/a_0) H \psi_0$:
  $c_s^2 \sim \Delta^2 \cdot a_0^2/c^2 = (c H \psi_0 / a_0)^2
  \cdot a_0^2/c^2 = H^2 \psi_0^2$.
  So $c_s \sim H \psi_0 \sim H \times 0.27$.

  The Jeans wavenumber is:
  $k_J = \sqrt{4\pi G \bar\rho} / c_s \sim H\sqrt{\Omega_m} / (H\psi_0)
  = \sqrt{\Omega_m}/\psi_0 \sim 2\;\text{Mpc}^{-1}$ (comoving).

  \textbf{This is a HUGE Jeans scale.} At $z \sim 1000$: $k_J \sim
  2$~Mpc$^{-1}$ (comoving), meaning modes with $k > 2$~Mpc$^{-1}$
  are supported against collapse. The BAO scale is $k_{\rm BAO}
  \sim 0.06$~Mpc$^{-1}$, safely below $k_J$. The nonlinear
  scale $k_{\rm NL} \sim 0.2$~Mpc$^{-1}$ is also below $k_J$.

  \textbf{BUT}: This $c_s$ applies to the psi-field perturbation
  equation with the \emph{background} $\Delta_{\rm bg}$ in the
  Newtonian regime. The dust-branch theorem applies in the
  opposite limit ($\Delta \to 0$). These are not contradictory
  --- they describe different regimes:

  \textbf{(a) Background evolution}: $\Delta_{\rm bg} \propto H(z)$,
  which is large at early times and $\sim 1.5$ today.
  The background psi-field is in the Newtonian (radiation-like) regime
  at all cosmologically relevant epochs.

  \textbf{(b) Perturbation growth}: Perturbations around this
  background feel an effective sound speed set by $K''(\Delta_{\rm bg})$.
  When $\Delta_{\rm bg} \gg 1$, $c_s \sim H\psi_0 \neq 0$, so
  perturbations DO have a Jeans scale.

  \textbf{(c) Late-time dust branch}: Only as $\Delta_{\rm bg} \to 0$
  (far future, $H \to 0$) does $c_s \to 0$ exactly.

  \textbf{CONCLUSION: The i16 claim that ``after crossover, psi-field
  perturbations cluster like CDM'' is WRONG. Throughout the entire
  observable history ($z < 3400$), the psi-field perturbation
  sound speed is $c_s \sim H\psi_0$, comparable to the Hubble
  expansion rate. This gives a Jeans length comparable to the Hubble
  radius. The psi-field CANNOT cluster on sub-Hubble scales at any
  epoch in the observable past.}

  This is the \textbf{existential crisis}.

\item \textbf{FINDING 3 (THE ESCAPE ROUTE --- $\Geff$ ENHANCEMENT):
  DFD does not need the psi-field to cluster independently.
  Baryonic matter clusters under $\Geff > G_N$, and the psi-field
  responds quasi-statically to the baryon distribution.
  The ``CDM-like clustering'' comes from the enhanced gravitational
  response, not from psi-field self-clustering.}

  The perturbation equation in DFD (section12\_cosmology.tex,
  Eq.~12.G) is:
  \begin{equation}
  \ddot\delta_k + 2H\dot\delta_k = 4\pi \Geff(k)\,\bar\rho_b\,\delta_k,
  \end{equation}
  where $\delta_k$ is the \emph{baryon} density contrast and
  \begin{equation}
  \Geff(k) = G_N \cdot \frac{k^2 + \kmu^2}{k^2}
  \end{equation}
  with $\kmu \approx 0.01$~Mpc$^{-1}$.

  At $k \gg \kmu$ (small scales, $\lambda \ll 600$~Mpc):
  $\Geff \approx G_N$ --- standard Newtonian growth.

  At $k \ll \kmu$ (large scales, $\lambda \gg 600$~Mpc):
  $\Geff \sim G_N \kmu^2/k^2 \gg G_N$ --- hugely enhanced growth.

  The critical question is: does the $\Geff$-enhanced growth of
  \emph{baryons} (which ARE pressureless dust after recombination)
  reproduce the CDM-like matter power spectrum?

  In \LCDM, the matter power spectrum at $z = 0$ is approximately:
  \begin{equation}
  P^{\rm CDM}(k) \propto k^{n_s} T^2(k),
  \end{equation}
  where $T(k)$ is the CDM transfer function with the characteristic
  turnover at $k_{\rm eq} \sim 0.01$~Mpc$^{-1}$ (set by the horizon
  at matter-radiation equality).

  In DFD, baryons are the ONLY clustering component. Before
  recombination, baryons are tightly coupled to photons (Silk
  damping, BAO). After recombination ($z < 1100$), baryons fall
  into the $\psi$-potential wells with $\Geff$-enhanced gravity.

  \textbf{The problem}: In \LCDM, CDM perturbations grow from
  $z_{\rm eq} \sim 3400$ while baryons are still coupled to photons.
  After decoupling at $z \sim 1100$, baryons fall into the
  pre-existing CDM potential wells and catch up. In DFD, there are
  no pre-existing potential wells. The psi-field responds
  quasi-statically to the matter distribution but does not
  independently cluster.

  \textbf{The $\Geff$ enhancement provides a partial rescue}:
  after recombination, baryons grow as
  $\delta_b \propto a^{p(k)}$ where the growth exponent $p(k)$
  satisfies:
  \begin{equation}
  p(p-1) + \frac{3}{2}p = \frac{3}{2}\frac{\Geff(k)}{G_N}
  = \frac{3}{2}\frac{k^2 + \kmu^2}{k^2}.
  \end{equation}
  At $k = \kmu$: $\Geff/G_N = 2$, giving $p \approx 1.28$
  (vs $p = 1$ for CDM). At $k = 0.1\,\kmu$: $\Geff/G_N = 101$,
  giving $p \approx 12$ (extremely rapid growth).

  However, on scales $k \gg \kmu$ (which is where most of the
  observable power spectrum lies, $k > 0.01$~Mpc$^{-1}$):
  $\Geff \approx G_N$ and $p \approx 1$, identical to CDM.

  \textbf{The real issue is on scales $0.001 < k < 0.1$~Mpc$^{-1}$}:
  the turnover region. In \LCDM, the turnover is set by the CDM
  transfer function. In DFD, the ``turnover'' is set by the baryon
  transfer function modified by $\Geff$. These will differ because:
  (a) baryons have Silk damping (oscillatory features below
  $k_{\rm Silk} \sim 0.1$~Mpc$^{-1}$),
  (b) there are no pre-existing CDM wells for baryons to fall into,
  (c) $\Geff$ enhancement partially compensates for the missing CDM
  head start.

  \textbf{Quantitative estimate}: The CDM growth advantage from
  $z_{\rm eq}$ to $z_{\rm dec}$ is a factor
  $D_{\rm CDM}(z_{\rm dec})/D_{\rm CDM}(z_{\rm eq})
  \approx a_{\rm dec}/a_{\rm eq} = (1+z_{\rm eq})/(1+z_{\rm dec})
  \approx 3400/1100 \approx 3$.
  This factor-of-3 head start is what allows CDM to form potential
  wells that baryons then fall into. In DFD, baryons must compensate
  with $\Geff > G_N$. On scales near $\kmu$, the enhanced growth
  exponent ($p \sim 1.3$) over the time from $z_{\rm dec}$ to $z = 0$
  gives a compensating factor of $(1100)^{0.28} \approx 5$, which
  MORE than compensates for the factor-of-3 CDM head start.

  \textbf{But on small scales} ($k \gg \kmu$): $\Geff = G_N$, and
  baryons grow at the same rate as CDM --- but starting from a
  SMALLER amplitude (no CDM wells to fall into). The deficit is
  exactly the factor $\sim 3$ that CDM gained between $z_{\rm eq}$
  and $z_{\rm dec}$.

  \textbf{VERDICT}: DFD's matter power spectrum matches \LCDM{} at
  $k \lesssim \kmu$ (enhanced growth compensates) but is suppressed
  by $\sim (1/3)^2 \sim 0.1$ in power at $k \gg \kmu$ (small scales)
  relative to \LCDM. This is a factor $\sim 10\times$ deficit in
  $P(k)$ at $k > 0.1$~Mpc$^{-1}$, or equivalently $\sigma_8$ would
  be $\sim 3\times$ too low.

  \textbf{This is a SERIOUS tension but NOT immediately fatal},
  because the actual baryon-photon dynamics are more complex than
  this simple estimate. The mochi\_class Boltzmann code is needed
  for a definitive answer. Furthermore, the quasi-static psi-field
  response modifies the effective potential in which baryons evolve,
  which is not simply ``$\Geff$ applied to baryons in vacuum.''

\item \textbf{FINDING 4 (CMB POWER SPECTRUM): The CMB $C_\ell$
  spectrum is largely SAFE because it is set at $z \sim 1100$
  when the psi-field is deep in the Newtonian regime.
  The psi-field acts as an effective smooth dark matter component
  for the CMB, entering only through $\Geff$ modification of the
  Sachs-Wolfe and ISW terms.}

  At recombination ($z \sim 1100$): $\Delta_{\rm bg} \sim 10^5 \gg 1$.
  The psi-field is deep in the Newtonian ($\mu \approx 1$) regime.
  The field equation reduces to:
  \begin{equation}
  \nabla^2\psi = -\frac{8\pi G}{c^2}\rho,
  \end{equation}
  i.e., standard Poisson. The psi-field tracks the matter distribution
  instantaneously (quasi-static response). The gravitational potentials
  seen by CMB photons at $z \sim 1100$ are therefore identical to
  \LCDM{} at leading order.

  The differences enter at:
  \begin{itemize}[nosep]
  \item \textbf{Integrated Sachs-Wolfe (ISW)}: At low $\ell$
    ($\ell < 30$), $\Geff > G_N$ modifies the potential decay rate,
    enhancing the ISW signal. This is the existing DFD prediction
    (Prediction 7, Cold Spot).
  \item \textbf{CMB lensing}: The lensing potential at $z \sim 0.5$--$2$
    is modified by $\Geff$, producing $A_{\rm lens} > 1$
    (existing Prediction 3).
  \item \textbf{Acoustic peaks}: At $z \sim 1100$, baryon-photon
    physics is standard. The psi-field does not participate in the
    acoustic oscillations (it has $c_s \sim H\psi_0$ which is much
    larger than the baryon-photon sound speed $c_s^{\rm b\gamma}
    \sim c/\sqrt{3}$ --- wait, $c_s^{\psi} \sim H\psi_0 \ll c/\sqrt{3}$
    at $z = 1100$: $H\psi_0 \sim 10^{-13}$~m/s at $z = 1100$
    vs $c/\sqrt{3} \sim 1.7\ee{8}$~m/s).
  \end{itemize}

  Actually, reconsidering: $c_s^{\psi} = H(z)\psi_0 \sim
  H_0(1+z)^{3/2}\sqrt{\Omega_m}\,\psi_0$. At $z = 1100$:
  $H \sim 10^5 H_0 \sim 2\ee{-13}$~s$^{-1}$, so
  $c_s^{\psi} \sim 2\ee{-13} \times 0.27 \sim 5\ee{-14}$~m/s.
  This is \textbf{extremely slow}. The psi-field Jeans length at
  $z = 1100$ is then:
  $\lambda_J = c_s / \sqrt{G\rho} \sim 5\ee{-14}/\sqrt{6.67\ee{-11}
  \times 10^{-18}} \sim 5\ee{-14}/2.6\ee{-14} \sim 2$~m (proper).
  This is \textbf{sub-meter}. The psi-field has essentially zero
  pressure support on all cosmological scales.

  \textbf{Wait --- this contradicts Finding 2.} Let me recompute.
  $c_s^2 = (W'/K'') \cdot (a_0^2/c^2)$.
  At $z = 1100$: $\Delta_{\rm bg} \sim 10^5$.
  $K'' = 1/(1+\Delta)^2 \approx 1/\Delta^2 = 10^{-10}$.
  $W' \approx 1$ (Newtonian regime for spatial sector).

  But what are the units? The perturbation equation (i16 P12/P15
  Eq.~in dust branch section) gives:
  $c_s^2 = W'(y_{\rm bg})/K''(\Delta_{\rm bg}) \cdot a_0^2/c^2$.
  $= 1/(10^{-10}) \times (1.2\ee{-10})^2/(3\ee{8})^2$
  $= 10^{10} \times 1.6\ee{-37}$
  $= 1.6\ee{-27}$~m$^2$/s$^2$.
  So $c_s \sim 4\ee{-14}$~m/s. This is consistent with the
  previous computation.

  Hmm. But this seems to say that the psi-field perturbations
  have $c_s \sim 10^{-14}$~m/s, i.e., they are essentially
  pressureless \textbf{at all times when $\Delta_{\rm bg} \gg 1$}.

  \textbf{RE-EXAMINATION}: The $c_s \sim c$ result from i16
  used $c_s^2 \sim H^2 d_H^2 / d_H^2 = H^2$ and then wrote
  ``$c_s \sim c$''. But $H^2$ has units of s$^{-2}$, not
  (m/s)$^2$. The i16 derivation confused $c_s^2$ (velocity squared)
  with an inverse-time-squared quantity. Let me redo carefully.

  From the perturbation equation:
  $K''(\Delta_{\rm bg})\,\ddot{\delta\psi}
  - W'(y_{\rm bg})\,\nabla^2\delta\psi = \text{source}$.

  For a plane wave $\delta\psi \propto e^{i(\mathbf{k}\cdot\mathbf{x}
  - \omega t)}$, the dispersion relation (ignoring source) is:
  $\omega^2 K''(\Delta_{\rm bg}) = k^2 W'(y_{\rm bg})$.

  But $k$ here is the physical wavenumber and $\omega$ is the physical
  frequency. The phase velocity is:
  $v_{\rm ph}^2 = \omega^2/k^2 = W'(y_{\rm bg})/K''(\Delta_{\rm bg})$.

  Now: what are the dimensions? $K$ and $W$ are both dimensionless
  functions appearing in the Lagrangian density
  $\mathcal{L} = (a_*^2/(8\pi G))[W(y) + K(\Delta)]$.
  The argument $y = |\nabla\psi|^2/a_*^2$ is dimensionless,
  $\Delta = (c/a_0)|\dot\psi - \dot\psi_0|$ is dimensionless.
  $K''$ and $W'$ are dimensionless.

  The kinetic term from $K$: $K(\Delta) \supset K''(\Delta)(\delta\Delta)^2
  /2$. With $\delta\Delta = (c/a_0)\delta\dot\psi$:
  kinetic energy $\propto K''(\Delta)(c/a_0)^2(\delta\dot\psi)^2$.

  The gradient term from $W$: $W(y) \supset W'(y) \cdot 2(\nabla\bar\psi
  \cdot \nabla\delta\psi)/a_*^2$. For the quadratic piece:
  $W'(y) |\nabla\delta\psi|^2 / a_*^2$.

  The dispersion relation is then:
  $K''(\Delta_{\rm bg})(c/a_0)^2 \omega^2 = W'(y_{\rm bg})/a_*^2 \cdot k^2$.

  So:
  \begin{equation}
  \boxed{c_s^2 = \frac{\omega^2}{k^2} = \frac{W'(y_{\rm bg})}{K''(\Delta_{\rm bg})}
  \cdot \frac{a_0^2}{a_*^2 c^2}}
  \end{equation}

  Since $a_* = 2a_0/c^2$: $a_*^2 = 4a_0^2/c^4$, so
  $a_0^2/(a_*^2 c^2) = a_0^2 c^4/(4a_0^2 c^2) = c^2/4$.

  Therefore:
  \begin{equation}
  c_s^2 = \frac{W'}{K''} \cdot \frac{c^2}{4}.
  \end{equation}

  At early times ($\Delta_{\rm bg} \gg 1$, $y_{\rm bg} \gg 1$):
  $W' \approx 1$, $K'' \approx 1/\Delta^2$:
  \begin{equation}
  c_s^2 \approx \Delta^2 \cdot \frac{c^2}{4}
  = \left(\frac{c}{a_0}H\psi_0\right)^2 \cdot \frac{c^2}{4}
  = \frac{c^4 H^2 \psi_0^2}{4 a_0^2}.
  \end{equation}
  At $z = 1100$: $H \sim 3.4\ee{5} H_0 \sim 7.5\ee{-13}$~s$^{-1}$:
  $c_s^2 \approx (3\ee{8})^4 \times (7.5\ee{-13})^2 \times 0.073
  / (4 \times (1.2\ee{-10})^2)$
  $= 8.1\ee{33} \times 5.6\ee{-25} \times 0.073 / (5.76\ee{-20})$
  $= 3.3\ee{8} / 5.76\ee{-20} = 5.7\ee{27}$~m$^2$/s$^2$.
  So $c_s \approx 2.4\ee{13}$~m/s $\gg c$.

  \textbf{THIS IS SUPERLUMINAL.} This means the perturbation equation
  is ill-posed (or rather, the perturbations are fully supported
  against collapse on all sub-Hubble scales). The psi-field
  perturbations cannot grow at $z = 1100$.

  \textbf{At late times} ($\Delta_{\rm bg} \to 1$):
  $K''(1) = 1/(1+1)^2 = 0.25$, $W'(y_0) \approx 1$:
  $c_s^2 \approx (c^2/4)/0.25 = c^2$.
  Still $c_s \sim c$.

  \textbf{As $\Delta_{\rm bg} \to 0$ (far future)}:
  $K''(0) = 1$, so $c_s^2 \to c^2/4$. STILL order $c$.

  \textbf{CRITICAL REALIZATION}: $c_s \to 0$ as $\Delta \to 0$
  only in the dust-branch proof where $p \to 0$ and $\rho \to \text{const}$.
  This describes the \emph{background} fluid's EoS, not the perturbation
  sound speed. The perturbation sound speed is $c_s^2 = dp/d\rho$
  evaluated from the thermodynamic relation, which for this k-essence
  system gives $c^2/4$ even in the dust limit.

  \textbf{No wait.} The Appendix Q proof says $c_s^2 \to 0$.
  Let me re-read: ``$dp/d\rho$ satisfies $dp/d\Delta = O(\Delta)$
  and $d\rho/d\Delta = \text{const} + O(\Delta)$, hence
  $c_s^2 = (dp/d\Delta)/(d\rho/d\Delta) \to 0$.''
  This is the adiabatic sound speed computed from the background
  thermodynamics, NOT from the perturbation dispersion relation.

  For k-essence, these two definitions can differ. The perturbation
  propagation speed is set by the kinetic matrix of the action,
  which I computed above as $c_s^2 = (W'/K'') \cdot c^2/4$.
  This does NOT go to zero.

  \textbf{THIS IS A FUNDAMENTAL INCONSISTENCY IN THE DFD DUST BRANCH
  CLAIM.} The background EoS has $w \to 0$, but the perturbation
  propagation speed is $c_s \sim c/2$ at all times. This means
  the psi-field is a stiff fluid (background dust, perturbation radiation)
  --- which is NOT CDM-like behavior for structure formation.

  \textbf{CMB impact}: The CMB is safe if and only if the psi-field
  does not participate in the acoustic oscillations as an independent
  degree of freedom. Since $c_s^{\psi} \sim c$, the psi-field
  perturbations propagate at the speed of light and are smoothed
  out on sub-Hubble scales. The psi-field contributes to the
  gravitational potential only through its quasi-static response
  to the matter distribution (via the $\Geff$ modified Poisson
  equation). This is qualitatively similar to a ``quintessence''
  dark energy with $c_s = c$, which is known to not cluster on
  sub-Hubble scales and to produce a standard CMB spectrum.

  \textbf{VERDICT}: The CMB $C_\ell$ spectrum survives because
  the psi-field is smooth at $z \sim 1100$. The matter power
  spectrum is the problem.

\item \textbf{FINDING 5 (DUST-BRANCH EOS: $w \sim 5\ee{-6}$ IS
  CORRECT; $10^{-37}$ WAS WRONG. BUT THE PERTURBATION SOUND SPEED
  IS $c_s \sim c/2$, NOT $c_s \to 0$. THESE ARE DIFFERENT QUANTITIES.}

  The i16 cpsi audit correctly identified that the present-day EoS
  is $w \sim \Delta_0 \sim 5\ee{-6}$ (not $10^{-37}$).
  The derivation:

  From Appendix Q: $p = (a_*^2/(8\pi G)) K(\Delta)$,
  $\rho = (a_*^2/(8\pi G))[(c/a_*)\dot\psi_0\,\Delta + O(\Delta^2)]$.
  For $\Delta \ll 1$: $K(\Delta) \approx \Delta^2/2$, so
  $w = p/\rho = (\Delta^2/2) / [(c/a_*)\dot\psi_0\,\Delta]
  = \Delta a_*/(2c\dot\psi_0)$.
  At present: $\Delta_0 \approx 1.5$ (from Finding 1), which gives
  $w_0 \approx 1.5/(2 \times c\dot\psi_0/a_*) \sim O(1)$.

  But wait: $\Delta_0 \approx 1.5$ is NOT in the $\Delta \ll 1$
  regime. The dust-branch expansion is only valid for $\Delta \ll 1$,
  which the present universe does NOT satisfy.

  For $\Delta = 1.5$: $K(1.5) = 1.5 - \ln(2.5) = 1.5 - 0.916 = 0.584$.
  $K'(1.5) = \mu(1.5) = 1.5/2.5 = 0.6$.
  $2\Delta K' - K = 2(1.5)(0.6) - 0.584 = 1.8 - 0.584 = 1.216$.
  $w = K/(2\Delta K' - K) = 0.584/1.216 = 0.48$.

  \textbf{$w \approx 0.48$ at present.} This is NOT dust. This is
  nearly radiation-like ($w = 1/2$ is the no-go lemma result for
  the wrong invariant).

  \textbf{HOWEVER}: This computation assumed the psi-field energy
  density is dominated by the temporal kinetic term. In reality,
  the total psi-field stress-energy includes both spatial ($W$)
  and temporal ($K$) sectors, plus the matter coupling.
  The full EoS depends on the relative magnitudes of these
  contributions.

  If the spatial gradient energy dominates (as expected for the
  background $\psi$-screen, where $|\nabla\psi|^2/a_*^2$ is the
  relevant parameter), then the temporal sector is subdominant
  and the ``dust branch'' result is only a statement about the
  temporal sector's contribution.

  \textbf{THE RESOLUTION}: In the DFD cosmological framework,
  the $\psi$-field is NOT treated as a separate fluid with its
  own $w$ and $c_s$. Instead, it enters through the modified
  Poisson equation ($\Geff$). The ``dust branch'' theorem
  establishes that the temporal sector does not introduce
  additional pressure that would prevent the theory from
  having a viable cosmological limit. The actual structure
  formation is governed by baryons growing under $\Geff$,
  not by psi-field self-clustering.

  \textbf{EoS derivation summary}:
  \begin{itemize}[nosep]
  \item i15 value $w \sim 10^{-37}$: \textbf{WRONG} (used $c_s^2/c^2$
    with wrong $c_s$)
  \item i16 cpsi audit value $w \sim 5\ee{-6}$: \textbf{WRONG}
    (used $\Delta_0 \ll 1$ expansion, but $\Delta_0 \approx 1.5$)
  \item Correct $w$ from full $K(\Delta)$ at $\Delta_0 = 1.5$:
    $w \approx 0.48$ for the temporal sector alone
  \item \textbf{Physical interpretation}: The temporal sector
    contributes a radiation-like component. Whether this is
    problematic depends on its energy density relative to
    the total matter/energy budget.
  \end{itemize}

  From the Appendix Q conserved current: $a^3 \mu(\Delta) = \text{const}$.
  Today: $a = 1$, $\mu(1.5) = 0.6$. At $z = 1100$: $a = 1/1101$,
  so $\mu(\Delta_{1100}) = 0.6 \times 1101^3 = 8\ee{8}$.
  But $\mu < 1$ always, so $\mu(\Delta_{1100}) \leq 1$.
  \textbf{Contradiction}: $a^3\mu(\Delta)$ cannot be conserved if
  $\mu < 1$ and $a^3$ grows by $10^9$.

  This means the conserved current $a^3\mu(\Delta) = \text{const}$
  only applies in the $\Delta \ll 1$ regime. For $\Delta \gg 1$
  (early universe), $\mu(\Delta) \approx 1$, and the conservation
  law becomes $a^3 = \text{const}$, which is obviously wrong.

  \textbf{The resolution}: The full conservation law is
  $a^3 K'(\Delta) \cdot (c/a_0) \cdot \text{sgn}(\dot\psi - \dot\psi_0)
  = \text{const}$, where the sign and magnitude of
  $\dot\psi - \dot\psi_0$ must be tracked. In the Newtonian regime
  ($\Delta \gg 1$): $K'(\Delta) = \mu(\Delta) \approx 1 - 1/\Delta$,
  and $\Delta \propto H(z) \propto (1+z)^{3/2}$, so
  $a^3 K'(\Delta) \approx a^3(1 - a^{3/2}/\Delta_0) \approx a^3$.
  The conservation law is not $a^3\mu = \text{const}$ but rather
  involves the full $\Delta(a)$ evolution.

  \textbf{Bottom line}: The EoS of the psi-field temporal sector
  is NOT $w \sim 5\ee{-6}$ and NOT $w \sim 10^{-37}$. At the
  present epoch ($\Delta_0 \approx 1.5$), $w \approx 0.48$. In the
  past ($\Delta \gg 1$), $w \to 1/2$. In the far future
  ($\Delta \to 0$), $w \to 0$. The dust branch is a future
  asymptotic, not the present state.

\end{enumerate}
\end{tcolorbox}


%=============================================================================
%=============================================================================
\part{Finding 1: The Transition Redshift}
\label{part:zcross}
%=============================================================================
%=============================================================================

\section{The Background Temporal Invariant}

From Appendix Q (Definition 2.1):
\begin{equation}
\Delta = \frac{c}{a_0}|\dot\psi - \dot\psi_0|.
\end{equation}

For the cosmological background $\psi$-screen with amplitude $\psi_0 \approx 0.27$
(from the distance-bias framework), the time derivative is:
\begin{equation}
\dot\psi_0 \sim H(z)\,\psi_0
\end{equation}
where $H(z) = H_0\sqrt{\Omega_m(1+z)^3 + \Omega_\Lambda}$ is the Hubble parameter.

\textbf{Critical question}: What enters $\Delta_{\rm bg}$?

Reading Appendix Q's Definition~2.1 carefully: $\Delta = (c/a_0)|\dot\psi - \dot\psi_0|$.
For the homogeneous background: $\psi = \psi_0$, so $\dot\psi = \dot\psi_0$ and
$\Delta_{\rm bg} = 0$.

For perturbations $\psi = \psi_0 + \delta\psi$:
$\Delta_{\rm pert} = (c/a_0)|\delta\dot\psi|$.

\textbf{The i16 derivation was using $\Delta_{\rm bg}$
= $(c/a_0)|\dot\psi_0|$ WITHOUT the background subtraction.}
This is dimensionally consistent but contradicts the Appendix Q definition
which explicitly subtracts $\dot\psi_0$.

If the Appendix Q definition is taken literally, then:
\begin{itemize}[nosep]
\item The background has $\Delta = 0$ (pure dust, $w = 0$ exactly).
\item Perturbations have $\Delta = (c/a_0)|\delta\dot\psi| \ll 1$
  for linear perturbations, giving $w \to 0$, $c_s^2 \to 0$.
\item There is NO ``radiation-to-dust transition'' at $z \sim 50$--$100$.
  The dust branch applies at all epochs for the perturbation sector.
\end{itemize}

\textbf{But then what is $\Delta_{\rm bg} = (c/a_0)H\psi_0$ that i16 computed?}

The i16 computation treated $\dot\psi_0$ itself as if it were a deviation
from zero. This would be the case if the ``background'' from which
deviations are measured is $\psi = 0$ (flat spacetime), not $\psi = \psi_0$
(cosmological screen). The Appendix Q definition uses $\psi_0$ as the
background, making $\Delta_{\rm bg} = 0$ by construction.

\subsection{Which interpretation is correct?}

\begin{proposition}[Two interpretations of $\Delta$]
\label{prop:delta-interp}
\hfill
\begin{enumerate}
\item \textbf{Interpretation A (Appendix Q literal)}: Background is $\psi_0$.
  $\Delta = (c/a_0)|\dot\psi - \dot\psi_0|$. Background has $\Delta = 0$.
  Perturbations have $\Delta \ll 1$. Dust branch applies at all epochs.
  No transition.

\item \textbf{Interpretation B (i16)}: Background is $\psi = 0$ (flat space).
  $\Delta = (c/a_0)|\dot\psi_0|$. Background has $\Delta \sim (c/a_0)H\psi_0$.
  Transition at $\Delta = 1$, i.e., $z \sim 50$--$100$.
\end{enumerate}
\end{proposition}

Under Interpretation A:
\begin{itemize}[nosep]
\item The psi-field perturbations have $c_s \to 0$ at all epochs.
\item BUT: the perturbation propagation speed from the kinetic matrix
  is $c_s^2 = (W'/K'') \cdot c^2/4$ evaluated at the BACKGROUND
  value. If $K''$ is evaluated at $\Delta_{\rm pert} \ll 1$ (since
  perturbations are small), then $K''(0) = 1$ and $c_s^2 = c^2/4$.
  This is ORDER $c$.

  However, the perturbation equation linearized around the background
  uses $K''$ evaluated at the background $\Delta_{\rm bg}$, not at
  the perturbation value. Under Interpretation A, $\Delta_{\rm bg} = 0$
  and $K''(0) = 1$, giving $c_s = c/2$.
\item This means the perturbation sound speed is $c/2$ regardless of
  epoch. The psi-field NEVER clusters on sub-Hubble scales.
\end{itemize}

Under Interpretation B:
\begin{itemize}[nosep]
\item At early times: $\Delta_{\rm bg} \gg 1$, $K''(\Delta) \ll 1$,
  $c_s \gg c$ (superluminal --- ill-posed or requires UV completion).
\item At late times: $\Delta_{\rm bg} \to 0$, $K''(0) = 1$,
  $c_s = c/2$.
\item Transition at $z \sim 50$--$100$ is from superluminal to $c/2$.
  Psi-field still does not cluster.
\end{itemize}

\textbf{Either way, the psi-field perturbation sound speed is at least
$c/2$ at all observable epochs. The psi-field DOES NOT cluster on
sub-Hubble scales. The ``dust branch'' refers only to the background
equation of state ($w \to 0$), not to the perturbation dynamics.}


%=============================================================================
%=============================================================================
\part{Finding 2: Does $\Geff$-Enhanced Growth Compensate?}
\label{part:geff}
%=============================================================================
%=============================================================================

\section{The DFD Growth Equation}

Since the psi-field does not independently cluster, DFD's structure
formation relies entirely on baryons growing under $\Geff$-enhanced
gravity. The growth equation is:
\begin{equation}
\ddot\delta_b + 2H\dot\delta_b = 4\pi \Geff(k)\,\bar\rho_b\,\delta_b,
\end{equation}
where $\delta_b$ is the baryon density contrast and
\begin{equation}
\Geff(k) = G_N \cdot \frac{k^2 + \kmu^2}{k^2},
\qquad \kmu \approx 0.01\;\text{Mpc}^{-1}.
\end{equation}

\subsection{Pre-recombination ($z > 1100$)}

Baryons are tightly coupled to photons. The baryon-photon fluid
undergoes acoustic oscillations with sound speed $c_s \approx c/\sqrt{3}$.
The acoustic oscillations produce the BAO feature.

In \LCDM, CDM perturbations grow logarithmically during radiation
domination ($z > z_{\rm eq}$) and then linearly during matter domination
($z_{\rm eq} > z > z_{\rm dec}$). By $z_{\rm dec} \approx 1100$,
CDM perturbations have amplitude $\delta_{\rm CDM} \sim 10^{-2}$.

In DFD, there is no CDM. The psi-field is smooth (does not cluster).
The gravitational potential at $z_{\rm dec}$ is sourced by baryons
(which have acoustic oscillations) and the smooth psi-field background.

\subsection{Post-recombination ($z < 1100$)}

Baryons decouple from photons and begin to fall freely under gravity.
In \LCDM: baryons fall into pre-existing CDM potential wells,
rapidly reaching $\delta_b \approx \delta_{\rm CDM} \sim 10^{-2}$
within a few $e$-folds.

In DFD: baryons must grow from their post-decoupling amplitude
$\delta_b(z_{\rm dec}) \sim 10^{-5}$ (acoustic oscillation amplitude)
under $\Geff$-enhanced gravity. On small scales ($k \gg \kmu$):
$\Geff = G_N$, and the growth is identical to a baryon-only universe.

\subsection{The baryon-only growth problem}

In a baryon-only universe with $\Geff = G_N$: the growth factor from
$z = 1100$ to $z = 0$ is $D(0)/D(1100) \sim 1100$ (matter-dominated
growth). Starting from $\delta_b \sim 10^{-5}$:
$\delta_b(0) \sim 10^{-5} \times 1100 \sim 10^{-2}$.

In \LCDM: starting from $\delta_{\rm CDM} \sim 10^{-2}$ at $z = 1100$:
$\delta_{\rm CDM}(0) \sim 10^{-2} \times 1100 \sim 10$.
(Nonlinear regime reached at $z \sim 10$.)

\textbf{DFD is a factor $\sim 1000\times$ behind \LCDM{} in the matter
density contrast at $z = 0$ on small scales ($k \gg \kmu$).}
This means $P_{\rm DFD}(k) / P_{\rm CDM}(k) \sim 10^{-6}$ at
$k \gg 0.01$~Mpc$^{-1}$.

\subsection{Can $\Geff$ rescue this?}

On scales $k \lesssim \kmu \approx 0.01$~Mpc$^{-1}$: $\Geff \gg G_N$.
The enhanced growth rate gives $\delta \propto a^p$ with $p > 1$.
At $k = 0.001$~Mpc$^{-1}$: $\Geff/G_N = (10^{-6} + 10^{-4})/10^{-6}
= 101$, giving $p \approx 12$. Growth from $z = 1100$:
$\delta(0) \sim 10^{-5} \times 1100^{12} \sim 10^{-5} \times 10^{36}
\sim 10^{31}$. This is vastly more than needed.

But on scales $k = 0.1$~Mpc$^{-1}$ (galaxy clusters):
$\Geff/G_N = (0.01 + 10^{-4})/0.01 = 1.01$. Growth:
$\delta(0) \sim 10^{-5} \times 1100^{1.005} \sim 10^{-5} \times 1103
\sim 1.1\ee{-2}$. Still $\sim 1000\times$ below CDM.

On scales $k = 1$~Mpc$^{-1}$ (galaxies): $\Geff/G_N \approx 1.0001$.
Completely negligible enhancement.

\begin{finding}[$\Geff$ cannot rescue small-scale power]
\label{find:geff-rescue}
The $\Geff$ enhancement operates only on scales $k \lesssim \kmu
\approx 0.01$~Mpc$^{-1}$ ($\lambda \gtrsim 600$~Mpc). On the
scales where the matter power spectrum is observed ($0.01 < k < 1$~Mpc$^{-1}$),
$\Geff \approx G_N$ and the growth is identical to a baryon-only
universe. DFD predicts a matter power spectrum that is
$\sim 10^{-6}\times$ lower than observed on scales $k > 0.1$~Mpc$^{-1}$.

\textbf{This is a catastrophic deficit.} The observed galaxy power
spectrum, cluster mass function, weak lensing shear, Lyman-$\alpha$
forest, and all small-scale structure formation probes would
disagree with DFD by many orders of magnitude.
\end{finding}


%=============================================================================
%=============================================================================
\part{Finding 3: The $P(k)$ Comparison and Escape Routes}
\label{part:Pk}
%=============================================================================
%=============================================================================

\section{Qualitative $P(k)$ Shape}

\begin{center}
\begin{tabular}{lccc}
\toprule
Scale & $k$ (Mpc$^{-1}$) & $P^{\rm DFD}/P^{\rm CDM}$ & Comment \\
\midrule
Hubble & $\sim 10^{-4}$ & $\gg 1$ & $\Geff$ vastly enhances \\
BAO turnover & $\sim 0.01$ & $\sim 1$ & $\Geff/G_N \sim 2$ \\
Cluster & $\sim 0.1$ & $\sim 10^{-3}$ & Baryon deficit \\
Galaxy & $\sim 1$ & $\sim 10^{-6}$ & Catastrophic deficit \\
Lyman-$\alpha$ & $\sim 10$ & $\sim 10^{-6}$ & Same \\
\bottomrule
\end{tabular}
\end{center}

\section{Escape Routes}

\subsection{Escape Route 1: Psi-field quasi-static response provides potential wells}

Even though the psi-field does not independently cluster (has $c_s \sim c$),
it responds quasi-statically to the matter distribution through the
modified Poisson equation:
\begin{equation}
\nabla\cdot[\mu(|\nabla\psi|/a_*)\nabla\psi] = -\frac{8\pi G}{c^2}\rho_b.
\end{equation}

This means the gravitational potential is not set by $\rho_b$ alone but
by $\rho_b / \mu$. For $\mu < 1$ (MOND regime), the potential is
\emph{deeper} than Newtonian for a given baryon distribution.

\textbf{But this IS the $\Geff$ effect already accounted for.} The
quasi-static response gives $\Geff(k) = G_N(k^2 + \kmu^2)/k^2$,
which we showed is insufficient on small scales.

\subsection{Escape Route 2: Two-component decomposition}

The two-component decomposition $\psi = \psi_{\rm grav} + \delta\psi_{\rm EM}$
(established in i5) could in principle allow $\psi_{\rm grav}$ and
$\delta\psi_{\rm EM}$ to have different propagation speeds. If
$\psi_{\rm grav}$ clusters at $c_s \to 0$ while $\delta\psi_{\rm EM}$
inherits EM speed, then the gravitational sector could provide CDM-like
potential wells.

\textbf{However}, the two-component decomposition is a mathematical
split of the same field $\psi$. Both components satisfy the same
field equation with the same $K(\Delta)$ temporal sector. There is
no mechanism for one component to have a different sound speed than
the other. The propagation speed is set by the kinetic structure of
the action, not by the source.

\textbf{Verdict}: The two-component decomposition does NOT provide
an escape route.

\subsection{Escape Route 3: Nonlinear growth from MOND dynamics}

In the nonlinear regime ($\delta > 1$), DFD's MOND-like dynamics
produce $a \sim \sqrt{a_N a_0}$ for $a < a_0$. This gives faster
collapse than Newtonian gravity for low-acceleration systems, potentially
accelerating structure formation once initial overdensities reach
the nonlinear regime.

\textbf{Problem}: The issue is reaching the nonlinear regime in the
first place. With $\delta_b(z=0) \sim 10^{-2}$ on galaxy scales,
the perturbations never go nonlinear, so the MOND enhancement
never activates.

\subsection{Escape Route 4: Modified initial conditions from inflation}

If the primordial power spectrum is enhanced relative to \LCDM{} on
small scales (blue tilt), the deficit could be partially compensated.
However, DFD does not predict a modified inflationary spectrum, and
Planck constrains $n_s = 0.965 \pm 0.004$ with no running.

\subsection{Escape Route 5: Psi-field IS CDM (reinterpretation)}

Perhaps the correct interpretation is that the psi-field's quasi-static
response IS the dark matter. In this view, there is no separate
``psi-fluid'' that needs to cluster; instead, the modified Poisson
equation $\Geff > G_N$ plays the role of CDM on all scales.

\textbf{The problem with this interpretation}: $\Geff/G_N = (k^2+\kmu^2)/k^2$
is only $> 2$ for $k < \kmu \approx 0.01$~Mpc$^{-1}$. On galaxy
scales ($k \sim 1$~Mpc$^{-1}$), $\Geff \approx G_N$. The extra
gravitational ``pull'' from the psi-field is negligible on small scales.
There is simply no mechanism to compensate for the missing CDM potential
wells between $z_{\rm eq}$ and $z_{\rm dec}$.

\subsection{Escape Route 6: Re-examine the perturbation sound speed}

The critical result is $c_s^2 = (W'/K'') \cdot c^2/4$. If either
$W'$ or $K''$ has a different form than assumed, $c_s$ could be
different. Specifically:

\textbf{If $K''$ is evaluated at the perturbation value $\Delta_{\rm pert}$
rather than the background $\Delta_{\rm bg}$}: For $\Delta_{\rm pert} \ll 1$,
$K''(\Delta_{\rm pert}) \approx 1$, giving $c_s = c/2$. No change.

\textbf{If the action normalization correction (i16 Finding 5) modifies
$K''$}: The i16 cpsi audit identified that Eq.~(6) of section02\_formalism
needs a $1/c^2$ correction factor. If this factor multiplies $K$, then
$K'' \to K''/c^2$ and $c_s^2 \to (W'/K'') \cdot c^4/4 \to$ even larger.
This makes things worse, not better.

\textbf{If the $a_0^2/(a_*^2 c^2) = c^2/4$ factor is wrong}:
Using $a_* = 2a_0/c^2$ gives $a_*^2 = 4a_0^2/c^4$ and
$a_0^2/(a_*^2 c^2) = c^2/4$. This is a direct consequence of the
$a_*$ definition and appears robust.

\begin{finding}[No viable escape route identified]
\label{find:no-escape}
All six escape routes fail. The fundamental problem is:
\begin{enumerate}[nosep]
\item The psi-field perturbation sound speed is $c_s \sim c/2$
  (from the kinetic matrix of the action), not $c_s \to 0$.
\item This prevents psi-field clustering on sub-Hubble scales.
\item Baryons alone cannot produce the observed matter power spectrum
  because they are tightly coupled to photons until $z \sim 1100$
  and lack the CDM ``head start'' from $z_{\rm eq}$ to $z_{\rm dec}$.
\item $\Geff$ enhancement operates only on scales $> 600$~Mpc,
  far larger than the scales where the deficit is catastrophic.
\end{enumerate}
\end{finding}


%=============================================================================
%=============================================================================
\part{Finding 4: CMB and Observational Implications}
\label{part:cmb}
%=============================================================================
%=============================================================================

\section{CMB Power Spectrum}

The CMB $C_\ell$ spectrum at $z \sim 1100$ is set by:
\begin{enumerate}[nosep]
\item Baryon-photon acoustic oscillations (peaks and troughs)
\item Silk damping (exponential suppression at high $\ell$)
\item Gravitational potential decay (ISW effect at low $\ell$)
\item CMB lensing (smearing of peaks at high $\ell$)
\end{enumerate}

In DFD:
\begin{itemize}[nosep]
\item Items 1--2 are identical to \LCDM{} if the psi-field is smooth
  at $z \sim 1100$ (confirmed: $c_s^{\psi} \sim c/2$).
\item Item 3 is modified by $\Geff$ (existing DFD prediction).
\item Item 4 is modified by $\Geff$ at low $z$ (existing DFD prediction).
\end{itemize}

\textbf{However}: The CMB acoustic peaks are sensitive to the
matter-radiation ratio $\omega_m = \Omega_m h^2$. In \LCDM,
$\omega_m$ includes CDM + baryons. In DFD, there is no CDM.
The ``effective $\omega_m$'' is set by the baryon density $\omega_b$
plus whatever contribution the smooth psi-field makes to the
expansion rate.

If the psi-field contributes to the background expansion as a fluid
with $w \sim 0.48$ (Finding 5), this is qualitatively different from
CDM ($w = 0$). The CMB peak structure would be modified.

\textbf{If the psi-field contributes to the expansion rate like CDM}
(which the distance-bias framework assumes --- DFD does not modify $H(z)$),
then the CMB peaks survive.

\textbf{The distance-bias framework}: DFD was established (i14) as a
theory that does NOT modify the Friedmann equation. $H(z)$ is identical
to \LCDM. The psi-screen biases distances but not expansion.

Under this framework: the background expansion has $\Omega_m = 0.315$
(which includes whatever the psi-field contributes). The psi-field's
contribution to $\rho$ must scale as $a^{-3}$ for the standard $H(z)$
to hold. From the conserved current: $a^3 \mu(\Delta) = \text{const}$,
so $\rho_\psi \propto a^{-3}$ if $\mu(\Delta) \propto \Delta$ (small
$\Delta$ regime). But at early times $\Delta \gg 1$ and
$\mu \approx 1$, so $\rho_\psi \propto a^{-3}$ trivially (the
psi-field energy density dilutes like matter).

\textbf{This is consistent}: the psi-field background energy density
scales as $a^{-3}$, mimicking CDM at the background level. The CMB
peaks are safe.

\textbf{The problem is entirely at the perturbation level}: the
psi-field perturbations do not cluster, so the potential wells
that drive baryon infall after decoupling are absent.


%=============================================================================
%=============================================================================
\part{Finding 5: The Dust-Branch EoS Derivation}
\label{part:eos}
%=============================================================================
%=============================================================================

\section{Three Estimates and Which is Correct}

\subsection{i15 estimate: $w \sim 10^{-37}$}

Used $w \sim c_s^2/c^2$ with $c_s \sim 10^{-5}$~m/s (from a specific
$c_\psi$ estimate). This is WRONG because: (a) $w = c_s^2/c^2$ is only
valid for an ideal fluid, not k-essence; (b) the $c_\psi$ value was
regime-dependent; (c) the perturbation sound speed is not the same as
the background EoS.

\subsection{i16 cpsi audit estimate: $w \sim 5\ee{-6}$}

Used the Appendix Q small-$\Delta$ expansion: $w \sim \Delta_0 \sim 5\ee{-6}$.
This assumed $\Delta_0 \ll 1$, but the actual value $\Delta_0 \approx 1.5$
(from $\Delta_0 = (c/a_0)H_0\psi_0$) is NOT small. The small-$\Delta$
expansion is invalid.

\textbf{WAIT}: Let me recompute $\Delta_0$ carefully.
$\Delta_0 = (c/a_0)|\dot\psi_0|$. In Appendix Q, $\dot\psi_0$ is
the \emph{background screen evolution rate}. If $\dot\psi_0 \sim H_0 \psi_0
= 2.2\ee{-18} \times 0.27 = 5.9\ee{-19}$~s$^{-1}$:
$\Delta_0 = (3\ee{8}/1.2\ee{-10}) \times 5.9\ee{-19}
= 2.5\ee{18} \times 5.9\ee{-19} = 1.48$.

But wait: the Appendix Q definition is $\Delta = (c/a_0)|\dot\psi - \dot\psi_0|$,
not $(c/a_0)|\dot\psi_0|$. For the BACKGROUND: $\dot\psi = \dot\psi_0$,
so $\Delta_{\rm bg} = 0$.

The i16 cpsi audit computed $\Delta_0 \sim H_0 c/a_0 \sim 5\ee{-6}$
(without the $\psi_0$ factor). Let me check:
$(c/a_0)H_0 = (3\ee{8}/1.2\ee{-10}) \times 2.2\ee{-18}
= 2.5\ee{18} \times 2.2\ee{-18} = 5.5$.

This gives $\Delta_0 \sim 5.5$, not $5\ee{-6}$. The $5\ee{-6}$ value
must have used $a_0/c$ instead of $c/a_0$:
$(a_0/c)H_0 = (1.2\ee{-10}/3\ee{8}) \times 2.2\ee{-18}
= 4\ee{-19} \times 2.2\ee{-18} = 8.8\ee{-37}$.

That gives $10^{-37}$, which is the i15 value. The i16 cpsi audit
line 983 says ``$w \sim \Delta_0 \sim H_0 c/a_0 \sim 5\ee{-6}$''.
But $H_0 c / a_0 = 2.2\ee{-18} \times 3\ee{8} / 1.2\ee{-10}
= 6.6\ee{-10}/1.2\ee{-10} = 5.5$, not $5\ee{-6}$.

There must be a factor issue. Perhaps they used $a_0 = 2\sqrt{\alpha} c H_0$
to write $c/(a_0/H_0) = c H_0 / a_0 = 1/(2\sqrt{\alpha}) \approx 5.8$.
Or perhaps $\Delta = (c/a_0) \dot\psi$ with $\dot\psi = H_0 \psi_0$
and $\psi_0 \sim 10^{-6}$? That would give
$\Delta_0 = 5.5 \times 10^{-6} = 5.5\ee{-6}$.

\textbf{AH}: If $\psi_0$ in the dust-branch context refers to the
\emph{perturbation} amplitude $\delta\psi \sim 10^{-6}$ (not the
screen amplitude $\psi_0 \sim 0.27$), then
$\dot{\delta\psi} \sim H_0 \delta\psi \sim 2.2\ee{-24}$~s$^{-1}$,
and $\Delta_{\rm pert} = (c/a_0) \times 2.2\ee{-24}
= 2.5\ee{18} \times 2.2\ee{-24} = 5.5\ee{-6}$.

This makes $\Delta_{\rm pert} \sim 5\ee{-6} \ll 1$, which IS in
the dust regime. And $w \sim \Delta \sim 5\ee{-6}$.

\textbf{RESOLUTION}: The i16 cpsi audit's $w \sim 5\ee{-6}$ is
the EoS of the PERTURBATION sector (treating $\delta\psi$ as a
separate fluid with $\Delta_{\rm pert} \ll 1$). This is consistent
with the Appendix Q dust branch theorem. The background sector
has $\Delta_{\rm bg} = 0$ and is trivially dust ($w = 0$ exactly).

\subsection{Correct interpretation}

\begin{equation}
\boxed{w_{\rm bg} = 0 \quad (\text{exact, by definition of background subtraction})}
\end{equation}
\begin{equation}
\boxed{w_{\rm pert} \sim \Delta_{\rm pert} \sim 5\ee{-6} \quad (\text{for linear perturbations at } z = 0)}
\end{equation}

The perturbation EoS is indeed $\sim 5\ee{-6}$, far below observational
bounds. The i16 cpsi audit was correct in its numerical value,
but the physical interpretation needs care: this describes the
effective EoS of the psi-field perturbation fluid, not the
background.

\textbf{BUT}: This $w_{\rm pert}$ describes the thermodynamic EoS
of the perturbation fluid, which characterizes the background pressure.
It does NOT determine the perturbation propagation speed $c_s^2$.
For k-essence, $c_s^2 \neq w \cdot c^2$. The perturbation propagation
speed is set by the second derivatives of the Lagrangian:
\begin{equation}
c_s^2 = \frac{\partial p / \partial X}{(\partial \rho / \partial X)}
\end{equation}
where $X$ is the kinetic variable. This can be $O(c^2)$ even when
$w \ll 1$.


%=============================================================================
\section*{Summary and Impact Assessment}
%=============================================================================

\begin{tcolorbox}[colback=red!10!white, colframe=red!70!black,
  title=\textbf{EXISTENTIAL ASSESSMENT}]

\textbf{STATUS: SERIOUS but NOT FATAL --- pending resolution of
perturbation sound speed.}

The critical issue is the perturbation propagation speed $c_s^2$
derived from the kinetic matrix of the DFD action. Our analysis
finds $c_s^2 = (W'/K'') \cdot c^2/4$, which is $O(c^2)$ at
all epochs, preventing psi-field clustering.

\textbf{However}, this derivation made specific assumptions about the
dimensional structure of the temporal-spatial coupling that remain
uncertain due to the i16-identified action normalization discrepancy.
The factor of $c^2/4$ arose from $a_0^2/(a_*^2 c^2) = c^2/4$ using
$a_* = 2a_0/c^2$. If the correct relationship between $a_*$ and
$a_0$ differs (as the action normalization correction might imply),
$c_s$ could be different.

\textbf{Three scenarios}:
\begin{enumerate}
\item \textbf{If $c_s \sim c$ (our derivation)}: DFD's psi-field
  cannot cluster. The matter power spectrum is catastrophically
  wrong on $k > 0.01$~Mpc$^{-1}$. DFD is \textbf{falsified at the
  cosmological level} unless the mochi\_class code reveals a
  mechanism not captured by this analytic treatment.

\item \textbf{If $c_s \to 0$ (the Appendix Q claim for perturbations)}:
  The psi-field clusters like CDM from $z_{\rm eq}$ onward.
  The matter power spectrum is viable. The CMB is safe.
  This requires that the perturbation propagation speed is set
  by a different mechanism than the kinetic matrix (e.g., if the
  quasi-static limit applies and the temporal kinetic term is
  negligible for the spatial Poisson equation).

\item \textbf{If the quasi-static approximation is the correct
  framework}: The psi-field responds instantaneously to the matter
  distribution through the modified Poisson equation, and the
  temporal kinetic term is irrelevant for structure formation.
  In this case, the ``perturbation sound speed'' question is moot:
  the psi-field is not a dynamical fluid but a constrained field
  determined by the matter distribution at each instant. This
  would make DFD equivalent to AQUAL cosmology, which is known
  to have its own (different) problems with the matter power spectrum.
\end{enumerate}

\textbf{RECOMMENDATION}: The mochi\_class Boltzmann code is now
EXISTENTIALLY URGENT. The analytic treatment cannot definitively
resolve the interplay between the temporal kinetic term and the
spatial quasi-static response. The code must:
\begin{enumerate}[nosep]
\item Implement the full perturbation equations (temporal + spatial)
\item Compute $P(k)$ at $z = 0$ in DFD
\item Compare with \LCDM{} and observed $P(k)$
\item Determine whether the quasi-static approximation is valid
\end{enumerate}

Until mochi\_class is operational, the dust branch's viability for
reproducing CDM-like clustering remains an \textbf{open question
at the existential level}.
\end{tcolorbox}


%=============================================================================
\section*{Corrections to Previous Iterations}
%=============================================================================

\begin{table}[H]
\centering
\caption{Corrections to i16 findings}
\begin{tabular}{clp{8cm}}
\toprule
\# & Finding & Correction \\
\midrule
1 & i16 P12/P15 Finding 3 & ``$z_{\rm cross} \sim 50$--$100$ is
  close to $z_{\rm eq}$'' is \textbf{WRONG}. $z_{\rm eq} = 3400$
  is $34$--$68\times$ higher. The transition redshift depends on
  interpretation of $\Delta$ (Finding 1). \\
2 & i16 P12/P15 Finding 3 & ``After crossover, psi-field clusters
  like CDM'' is \textbf{WRONG or at best UNPROVEN}. The perturbation
  propagation speed $c_s$ from the kinetic matrix is $O(c)$ at all
  epochs (Finding 2). \\
3 & i16 ITERATION\_TRACKER & ``Dust-branch EoS corrected: $w \sim 5\ee{-6}$
  (not $10^{-37}$)'' is \textbf{CORRECT in value but the derivation
  pathway was unclear}. $w \sim 5\ee{-6}$ applies to the perturbation
  sector with $\Delta_{\rm pert} \ll 1$. The background has $w = 0$
  exactly (Finding 5). \\
4 & i16 P12/P15 Finding 3 & ``Naturally reproducing the CDM transfer
  function'' is \textbf{UNPROVEN and likely WRONG} without the
  mochi\_class code. \\
\bottomrule
\end{tabular}
\end{table}


%=============================================================================
\section*{Open Questions for i18}
%=============================================================================

\begin{enumerate}
\item \textbf{PRIORITY 1}: Is the quasi-static approximation valid for
  DFD perturbations on sub-Hubble scales? If so, the temporal kinetic
  term is irrelevant and the perturbation dynamics are set by the
  modified Poisson equation. This must be verified by comparing
  $\omega_K = \sqrt{k^2 c_s^2}$ (kinetic timescale) with $H$
  (expansion timescale) and $\omega_{\rm grav} = \sqrt{4\pi G\rho}$
  (gravitational timescale).

\item \textbf{PRIORITY 2}: In AQUAL cosmology (the closest analog to
  DFD), what is the status of the matter power spectrum? Skordis and
  Z\l{}o\'snik (2021) showed that their theory (RMOND) can match
  $P(k)$ by adding a ``dark fluid'' with specific properties. Does
  DFD's temporal sector play the role of this dark fluid?

\item \textbf{PRIORITY 3}: Does the action normalization correction
  (i16 Finding 5) change $c_s$?

\item \textbf{PRIORITY 4}: The psi-field perturbation equation should
  be derived from first principles using the FULL action (spatial +
  temporal + matter coupling). The i16 P12/P15 derivation used only
  the temporal and spatial kinetic terms without the source term.

\item \textbf{PRIORITY 5}: Skordis-Z\l{}o\'snik RMOND achieved
  CDM-like $P(k)$ by having a dark field with TWO kinetic terms
  (standard + ``dark force''). Does DFD's two-sector structure
  (W + K) provide an analogous mechanism?
\end{enumerate}

\vspace{2em}

\textit{Filed 20 March 2026. Agent 12 (Cosmology --- Dust Branch
+ Late Clustering), Opus 4.6, Wave 4 Iteration 17.}

\textit{New results: 5 (transition redshift re-derivation,
perturbation sound speed crisis, $\Geff$-rescue failure,
CMB safety analysis, EoS correction hierarchy).
Corrections to i16: 4 (transition redshift misidentified,
clustering claim unproven, EoS derivation pathway, transfer function
claim). Status: EXISTENTIAL QUESTION OPEN. mochi\_class now
critical path.}

\end{document}
